\section{The `Maintenance Trio': Glymphatic, Synaptic, and Cellular}\label{sec:maintenance_trio}

Before the brain can process information, it must be physically maintained. This housekeeping occurs primarily during NREM deep sleep (N3).

\begin{itemize}
    \item \textbf{Glymphatic Clearance (The Brain's ``Wash Cycle''):} During N3, the brain's glial cells are believed to shrink, increasing interstitial space by up to 60\%. This allows cerebrospinal fluid (CSF) to flush through the brain tissue, clearing metabolic byproducts that accumulate during wakefulness, such as amyloid-beta (a protein implicated in Alzheimer's disease) \citep{benington1999, xie2023}.
    
    \item \textbf{Synaptic Homeostasis (Pruning the Network):} During wakefulness, we are constantly forming new synaptic connections as we learn. This is energetically unsustainable. The \textit{Synaptic Homeostasis Hypothesis} (SHY) posits that NREM sleep performs a vital ``pruning'' function, weakening and removing weaker, non-essential synaptic connections. This process ``cleans the slate'', reducing the brain's energy cost and improving the signal-to-noise ratio for new learning the next day \citep{price_of_plasticity2014}.
    
    \item \textbf{Cellular Maintenance:} Sleep also supports cellular-level restoration and metabolic housekeeping, coordinating processes that favor repair, clearance, and remodeling.
\end{itemize}
